\subsubsection{\texttt{csapi}}

使用\texttt{csapi}函数函数时请至少输入四个插值点，更少插值点情形的功能待完善。

\texttt{csapi}函数根据给定的插值点列使用非扭结边界条件(not-a-knot)，共有2种用法:

\textbf{I.}  \texttt{pp = csapi(x,y)},该函数返回值为一个结构体(pp 格式)，其中

\begin{itemize}
    \item \texttt{x y} 分别为样条插值坐标点的横纵坐标构成的一维向量，长度必须相同,且要求\texttt{x}已经完成单增排列；
    \item \texttt{pp} 为一结构体，包含如下字段: \texttt{form} (样条的形式，在这里为 pp 格式)、\texttt{breaks} (样条节点的位置，为一维向量)、\texttt{coefs}(分片多项式系数组成的矩阵，每行为一个三次多项式的系数)、\texttt{pieces}(样条的多项式片数)、\texttt{order} (分片多项式的阶数)、\texttt{dim} (目标函数的维数);
\end{itemize}

如输入：
\begin{lstlisting}[language = matlab]
    x = [0 1 2 3]; 
    y = [1 2 4 6]; 
    s = csapi(x,y)
\end{lstlisting}

输出为：
\begin{lstlisting}[language = matlab ]
    form  : 'pp'
    breaks: [0 1 2 3]
    order : 4
    coefs : [3 * 4 double]
    pieces: 3
    dim   : 1
\end{lstlisting}

\textbf{II.}  \texttt{values = csapi(x,y,xx)},该函数返回生成样条函数在xx处的值，其中
\begin{itemize}
    \item \texttt{x y} 分别为样条插值坐标点的横纵坐标构成的一维向量，长度必须相同,且要求\texttt{x}已经完成单增排列；
    \item \texttt{xx} 为待求值点横坐标构成的一维向量;
    \item \texttt{value} 为一维向量，返回待求值点的函数值;
\end{itemize}

如输入：
\begin{lstlisting}[language = matlab]
    x = [0 1 2 3]; 
    y = [0 1 4 9]; 
    xx = [0.5 1.5]; 
    values = csapi(x,y,xx)
\end{lstlisting}

输出为：
\begin{lstlisting}[language = matlab ]
    values = [0.7500 2.7500]
\end{lstlisting}